\documentclass[11pt]{article}

\usepackage{amsmath, amssymb, bm, mathtools}
\usepackage{hyperref}
\usepackage{graphicx}
\usepackage{color}
\usepackage{xeCJK}

\title{Applied Linear Regression \\ 应用回归分析}
\author{蒋智超 \thanks{UC Berkeley, 2025 Fall} \\
        \normalsize 教材: Linear Model and Extensions (Peng Ding, arxiv:2401.00649)}
\date{\today}

\begin{document}

\maketitle

\tableofcontents

\section{课程简介}

本课程使用 Peng Ding 的《Linear Model and Extensions》作为教材，涵盖以下主题：

\begin{itemize}
\item 最小二乘法 (OLS) 及其性质
\item Gauss-Markov 定理与 BLUE
\item 正态线性模型与统计推断
\item 正则化方法：岭回归与 Lasso
\item 广义线性模型 (GLM)
\item Logistic 回归
\item 计数数据模型
\end{itemize}

\section{课件生成进度}

\subsection{Gauss-Markov (Chapter 04)}
状态: \textbf{Pending}

\subsection{Normal Linear Model (Chapter 05)}
状态: \textbf{Pending}

\subsection{OLS Asymptotics (Chapter 03)}
状态: \textbf{Pending}

\section{参考文献}

\bibliographystyle{plain}
\bibliography{../PDFs/applied-linear-regression/linearmodel_lecturenotes_pengding}

\end{document}
